\section{A conjectural refinement forced by Fourier symmetry}
\label{sec:signed-orbit-conjecture}

The relation $F_E^2|x\rangle=|-x\rangle$ means that retaining arithmetic
Fourier also retains negation. Negation commutes with Frobenius, but in
odd characteristic it does not fix $1$ and is not a field automorphism.
This is a closure of arithmetic operators, distinct from the normal
closure of a field extension. It suggests a concrete test of how much
additional information a filtered boundary must carry.

\begin{definition}[Signed Frobenius orbit candidates]
Let $p$ be odd and $r\geq1$. The generators of
$\Gamma_r^{\rm sgn}=C_r\times C_2$ act on
$\mathbb F_{p^r}^{\times}$ by $x\mapsto x^p$ and $x\mapsto-x$.
A point of exact Frobenius period $d$ is internal when its negative lies
in its Frobenius orbit, and external otherwise. Let $t$ be a real
variable. The integer M\"obius function $\mu$ is zero on integers divisible
by a prime square and otherwise $(-1)$ to the number of prime factors.
Write $\varphi(d)=|\{1\leq a\leq d:\gcd(a,d)=1\}|$ and
$c_d(t)=\sum_{e\mid d}\mu(d/e)t^e$.
For $d=2^k m$ even, with $k\geq1$ and $m$ odd, define
\[
 A_d^{\rm sgn}(t)=\sum_{e\mid m}\mu(m/e)(t^{2^{k-1}e}-1),\qquad
 B_d^{\rm sgn}(t)=\sum_{e\mid m}\mu(m/e)(t^{2^{k-1}e}-1)^2.
\]
For odd $d$ set $A_d^{\rm sgn}=0$, $B_1^{\rm sgn}=t-1$, and
$B_d^{\rm sgn}=c_d(t)$ for $d>1$. Finally put
\[
 J_2(m)=m^2\prod_{\substack{\ell\mid m\\\ell\text{ prime}}}
                      (1-\ell^{-2}),
\]
with empty product one.
\end{definition}
\begin{scope}
These are candidate count polynomials. Characteristic two is excluded:
there negation is identity and the group action changes.
\end{scope}
\provenance{D1331, D1441}{definitions.md}
{theory/checks/arithmetic\_signed\_check.py}{conjectural increment}

\begin{conjecture}[First- and second-order signed orbit sectors]
\statusConjecture\quad
For every odd prime $p$, every $r\geq1$, and every $d\mid r$,
$A_d^{\rm sgn}(p)$ and $B_d^{\rm sgn}(p)$ count the internal and
external nonzero labels of exact Frobenius period $d$, respectively.
Internal signed orbits occur only for even $d$ and have size $d$;
external signed orbits have size $2d$. Every nonzero candidate polynomial
is positive for all $t>1$. For even $d=2^k m$,
\begin{align*}
 A_d^{\rm sgn}(t)&=\varphi(d)(t-1)+O((t-1)^2),\\
 B_d^{\rm sgn}(t)&=2^{2k-2}J_2(m)(t-1)^2+O((t-1)^3).
\end{align*}
For odd $d$, including $d=1$, the first coefficient of
$B_d^{\rm sgn}(t)$ is $\varphi(d)$.
\end{conjecture}
\begin{scope}
The quantifiers range over all degrees and all odd primes. The finite
census supports this conjecture and does not prove it. No full Fourier,
multiplication or mixed-process limit follows from the proposed counts.
\end{scope}
\provenance{LIM-SIGNED}{docs/research-plans/signed-orbit-filtration.md}
{theory/checks/arithmetic\_signed\_check.py, S1--S2}
{not submitted for theorem admission in this round}

For a degree-two extension, the formula predicts the following three
types of nonzero labels.
\begin{center}
\begin{tabular}{llll}
\toprule
Type & Point count & Signed orbit size & Vanishing order\\
\midrule
Prime-field labels & $p-1$ & $2$ & $1$\\
Internal period-two labels & $p-1$ & $2$ & $1$\\
External period-two labels & $(p-1)^2$ & $4$ & $2$\\
\bottomrule
\end{tabular}
\end{center}
The two first-order types suggest two matrix blocks $M_2$: on the first,
Frobenius is identity and negation exchanges the points; on the second,
both exchange them. The second-order type suggests $M_4$ with the regular
$C_2\times C_2$ action. This is a specific proposed mechanism by which
closing the arithmetic symmetries retains additional orders of degeneration.
It does not identify those orders with arbitrary Clifford levels.

The exact checker enumerates quotient fields of orders $3,9,27,81,5,25$,
obtains periods by iteration, and determines internality by orbit membership
of the actual negative. It separately expands the sparse polynomials about
one through degree $48$. Both the incorrect-negation and incorrect-order
mutations fail. The conjecture's general proof and its integration with
Fourier remain future work.
